\documentclass{article}

% ── Core packages ─────────────────────────────────────────────────────────────
\usepackage[utf8]{inputenc}
\usepackage[T1]{fontenc}
\usepackage[numbers]{natbib}
\DeclareUnicodeCharacter{2014}{---}
\DeclareUnicodeCharacter{2192}{$\to$}
\DeclareUnicodeCharacter{00D7}{$\times$}
\DeclareUnicodeCharacter{2605}{$\star$}
\usepackage{hyperref}
\usepackage{url}
\usepackage{booktabs}
\usepackage{amsfonts}
\usepackage{amsmath}
\usepackage{amssymb}
\usepackage{nicefrac}
\usepackage{microtype}
\usepackage{xcolor}
\usepackage{graphicx}
\usepackage{subcaption}
\usepackage{multirow}
\usepackage{array}
\usepackage{enumitem}
\usepackage{float}
\usepackage{wrapfig}
\usepackage{tcolorbox}
\usepackage{mdframed}
\usepackage{soul}       % for \hl{}
\usepackage{geometry}

% ── Page geometry (slightly wider for tables) ─────────────────────────────────
\geometry{left=1.5cm, right=1.5cm, top=2cm, bottom=2.5cm}

% ── Custom colours ────────────────────────────────────────────────────────────
\definecolor{mblue}{HTML}{2563EB}
\definecolor{mgreen}{HTML}{16A34A}
\definecolor{mred}{HTML}{DC2626}
\definecolor{morange}{HTML}{D97706}
\definecolor{mpurple}{HTML}{7C3AED}
\definecolor{mlightgray}{HTML}{F3F4F6}
\definecolor{mboxbg}{HTML}{EFF6FF}
\definecolor{mboxborder}{HTML}{2563EB}

% ── Insight box ───────────────────────────────────────────────────────────────
\newtcolorbox{insightbox}[1][]{
  colback=mboxbg,
  colframe=mboxborder,
  fonttitle=\bfseries\small,
  title={#1},
  boxrule=0.8pt,
  arc=3pt,
  left=6pt, right=6pt, top=4pt, bottom=4pt
}

\newtcolorbox{warningbox}[1][]{
  colback=yellow!10,
  colframe=morange,
  fonttitle=\bfseries\small,
  title={#1},
  boxrule=0.8pt,
  arc=3pt,
  left=6pt, right=6pt, top=4pt, bottom=4pt
}

% ── Hyper setup ───────────────────────────────────────────────────────────────
\hypersetup{
  colorlinks=true,
  linkcolor=mblue,
  citecolor=mgreen,
  urlcolor=mpurple
}

% Consistent figure formatting
\newcommand{\reportfigure}[3]{%
  \begin{figure}[tbp]
    \centering
    \includegraphics[width=\linewidth,height=0.78\textheight,keepaspectratio]{#1}
    \caption{#2 This figure shows #3}
  \end{figure}
}

% ─────────────────────────────────────────────────────────────────────────────
\title{%
  \textbf{Dissecting a Decade of Vietnamese E-Commerce Revenue}\\[4pt]
  {\large A Full Exploratory Data Analysis and Preprocessing Report}\\[2pt]
  {\normalsize Vin Datathon — Daily Revenue Forecasting Challenge}
}

\author{%
  Team Vin Datathon\\
  \texttt{vin-datathon-eda-v1.0}\\[2pt]
  \small\textit{Dataset: 2012-07-04 $\to$ 2022-12-31 $\;\cdot\;$
  Forecast horizon: 2023-01-01 $\to$ 2024-07-01 (548 days)}
}

\begin{document}
\maketitle

% ── Abstract ──────────────────────────────────────────────────────────────────
\begin{abstract}
We present a comprehensive exploratory data analysis (EDA) of a Vietnamese fashion
e-commerce enterprise covering \textbf{3,833 daily observations} across thirteen
relational tables — sales, customers, orders, products, promotions, shipments,
returns, reviews, payments, geography, inventory, and web traffic.
Our analysis reveals four structurally important phenomena that any forecast
model must capture: (i)~a \textbf{sharp regime break in 2018--2019} accompanied
by a $-39\%$ year-over-year revenue collapse, (ii)~a \textbf{remarkably stable
annual seasonal shape} ($\approx2.7\times$ peak-to-trough amplitude) that
persists across the regime break, (iii)~a \textbf{biennial August anomaly} in
which August revenue alternates between even and odd years with Pearson
$r = 0.97$, and (iv)~a pronounced \textbf{COD cancellation drag} that inflates
operational cost invisibly within reported revenue.
We additionally document the feature engineering pipeline — cyclical encodings,
leakage-safe long-horizon lags, and seasonal lookup tables — that translates
these findings into model-ready signals for the 548-day forecast window.
\end{abstract}

% ─────────────────────────────────────────────────────────────────────────────
\section{Introduction}
\label{sec:intro}
% ─────────────────────────────────────────────────────────────────────────────

Forecasting daily revenue in the context of Vietnamese fashion retail is
considerably harder than it may appear from the raw numbers.
The dataset spans a full decade (July 2012 to December 2022), but the business
underwent a structural shock around 2018--2019 that effectively renders the
pre-break years a different data-generating process.
The task — predicting revenue for every calendar day from January 2023 through
July 2024 — therefore demands a careful reconciliation of two questions:
\textit{which historical patterns are structural and will persist?} and
\textit{which patterns belong to a regime that no longer exists?}

This report is organised as follows.
Section~\ref{sec:data} describes the dataset and performs a data quality audit.
Section~\ref{sec:ts} analyses the revenue time series in depth: trend, seasonal
decomposition, autocorrelation structure, stationarity tests, and spectral
content.
Section~\ref{sec:customers} covers customer acquisition, purchase frequency, and
RFM segmentation.
Section~\ref{sec:products} examines the product and category landscape.
Section~\ref{sec:ops} investigates operational signals — promotions, shipping,
returns, reviews, web traffic, and payment behaviour.
Section~\ref{sec:fe} documents the preprocessing and feature engineering
decisions.
Section~\ref{sec:findings} synthesises eight actionable hypotheses.

\begin{insightbox}[Reading guide]
Throughout the report, blue boxes like this one contain analytical insights
that directly inform downstream modelling decisions.
Orange boxes flag data quality concerns that require careful handling.
\end{insightbox}

% ─────────────────────────────────────────────────────────────────────────────
\section{Dataset Description and Data Quality}
\label{sec:data}
% ─────────────────────────────────────────────────────────────────────────────

\subsection{Schema Overview}

The dataset comprises 13 relational tables.
Table~\ref{tab:schema} summarises their sizes, temporal coverage, and notable
missingness.

\begin{table}[H]
\centering
\caption{Schema overview — all 13 tables.
$^\dagger$Rows for \texttt{order\_items} are line-item records (not orders).}
\label{tab:schema}
\small
\begin{tabular}{lrrlll}
\toprule
\textbf{Table} & \textbf{Rows} & \textbf{Cols} & \textbf{Date range}
  & \textbf{Key identifier} & \textbf{Notable missing} \\
\midrule
\texttt{sales}       & 3,833   & 3  & 2012-07-04 → 2022-12-31 & Date          & None \\
\texttt{customers}   & $\sim$50k   & 8  & 2012-07 → 2022-12       & customer\_id  & None \\
\texttt{orders}      & $\sim$200k  & 9  & 2012-07 → 2022-12       & order\_id     & Few zips \\
\texttt{order\_items}& $\sim$600k$^\dagger$ & 8 & — & order\_id+product\_id & promo\_id ($\sim$60\%) \\
\texttt{products}    & $\sim$1k    & 6  & —                       & product\_id   & None \\
\texttt{promotions}  & $\sim$80    & 6  & 2012-07 → 2022-12       & promo\_id     & None \\
\texttt{shipments}   & $\sim$180k  & 5  & 2012-07 → 2022-12       & shipment\_id  & None \\
\texttt{returns}     & $\sim$12k   & 5  & 2012-07 → 2022-12       & return\_id    & None \\
\texttt{reviews}     & $\sim$90k   & 5  & 2012-07 → 2022-12       & review\_id    & None \\
\texttt{payments}    & $\sim$200k  & 5  & 2012-07 → 2022-12       & payment\_id   & None \\
\texttt{geography}   & $\sim$5k    & 5  & —                       & zip           & None \\
\texttt{inventory}   & $\sim$40k   & 7  & monthly snapshots        & product\_id   & None \\
\texttt{web\_traffic}& $\sim$3.8k  & 4  & 2012-07 → 2022-12       & date          & None \\
\bottomrule
\end{tabular}
\end{table}

\subsection{Sales Series Completeness}

The core \texttt{sales} table is the forecasting target.
It contains \textbf{exactly one record per calendar day} with no gaps across the
entire 3,833-day range, and zero revenue/COGS anomalies.
This is an unusually clean time series — a notable contrast to most real-world
retail datasets where weekends, holidays, and system outages create missing dates.

\subsection{Foreign Key Integrity}

A systematic foreign-key check across all relational joins reveals
$>99\%$ coverage in all critical paths:
\texttt{orders.customer\_id} $\to$ \texttt{customers},
\texttt{order\_items.order\_id} $\to$ \texttt{orders},
\texttt{order\_items.product\_id} $\to$ \texttt{products},
\texttt{shipments.order\_id} $\to$ \texttt{orders},
\texttt{returns.order\_id} $\to$ \texttt{orders}.
The only slightly lower coverage ($\sim$96\%) is
\texttt{orders.zip} $\to$ \texttt{geography}, reflecting a small number of
non-standard zip codes.

\subsection{Missingness Patterns}

The most notable missingness is in \texttt{order\_items.promo\_id}, where
approximately 60\% of line items carry a \texttt{NULL} value.
This is \emph{semantically correct} — it indicates that most items were purchased
without a promotion code, not that the data is corrupted.
Treating \texttt{NULL promo\_id} as ``no promotion'' is therefore the correct
encoding, giving a meaningful binary flag.

\begin{warningbox}[Data quality note — promo\_id nulls]
\texttt{promo\_id} is \texttt{NULL} for the majority of order-item lines.
\textbf{Do not impute} this field; instead, create a binary
\texttt{has\_promo} indicator.
The null itself is the signal: it means the item was sold at full price.
\end{warningbox}

% ─────────────────────────────────────────────────────────────────────────────
\section{Revenue Time-Series Analysis}
\label{sec:ts}
% ─────────────────────────────────────────────────────────────────────────────

\subsection{Annual and Monthly Overview}

Table~\ref{tab:annual} tells the macro story succinctly.
Revenue grew from approximately 0.90~B~VND in 2012 to a peak of
\textbf{2.10~B~VND in 2016}, then entered a multi-year decline culminating in
a \textbf{$-39\%$ collapse in 2019}.
The years 2019--2022 form a distinct ``recovery regime'' at a structurally lower
absolute level ($\sim$1.1--1.5~B~VND/year), with the most recent data point
(2022, $+14.8\%$) providing a cautiously optimistic signal for the forecast
horizon.

\begin{table}[H]
\centering
\caption{Annual revenue summary (Billion VND). Gross margin is consistently
high throughout, suggesting that the decline is demand-driven rather than
margin-squeeze.}
\label{tab:annual}
\small
\begin{tabular}{lrrrrr}
\toprule
\textbf{Year} & \textbf{Revenue (B)} & \textbf{COGS (B)}
  & \textbf{Gross Profit (B)} & \textbf{Margin \%} & \textbf{YoY \%} \\
\midrule
2012 & 0.90 & 0.77 & 0.13 & 14.2 & — \\
2013 & 1.12 & 0.96 & 0.16 & 14.3 & +24.4 \\
2014 & 1.43 & 1.22 & 0.21 & 14.7 & +27.7 \\
2015 & 1.78 & 1.52 & 0.26 & 14.6 & +24.5 \\
2016 & \textbf{2.10} & 1.79 & 0.31 & 14.8 & +18.0 \\
2017 & 1.96 & 1.67 & 0.29 & 14.8 & $-6.7$ \\
2018 & 1.87 & 1.59 & 0.28 & 15.0 & $-4.6$ \\
2019 & 1.14 & 0.97 & 0.17 & 14.9 & $-38.9$ \\
2020 & 1.18 & 1.00 & 0.18 & 15.2 & $+3.5$ \\
2021 & 1.03 & 0.88 & 0.15 & 14.6 & $-12.7$ \\
2022 & 1.18 & 1.00 & 0.18 & 15.3 & $+14.8$ \\
\bottomrule
\end{tabular}
\end{table}

\reportfigure{saved_01_monthly_revenue_margin.png}{Revenue trajectory and margin stability.}{the decade-long revenue trajectory, the stable gross-margin band, and the structural break around 2018--2019, and supports regime-aware modelling of revenue level versus margin behaviour.}

\begin{insightbox}[Key observation — the margin is surprisingly stable]
Gross margin hovers in the remarkably tight band of \textbf{14.2\%--15.3\%} for
the entire decade, across both the growth and the collapse.
This rules out price wars, cost inflation, or product mix changes as the cause
of the 2019 collapse.
The shock was almost certainly demand-side: a loss of customers or traffic,
not a deterioration of the unit economics.
\end{insightbox}

\subsection{The 2018--2019 Structural Break}

The $-39\%$ revenue drop from 2018 to 2019 is the single most important feature
of this dataset.
Its causes are not directly observable in the data, but several indirect signals
are consistent:
\begin{itemize}[leftmargin=1.5em, itemsep=2pt]
  \item The drop occurs \emph{instantaneously} at the year boundary — January
    2019 is dramatically lower than December 2018 — indicating a sudden event
    rather than a gradual trend.
  \item New customer acquisition (Section~\ref{sec:customers}) shows a sharp
    deceleration starting in mid-2018.
  \item Promo penetration jumps in 2019--2021 (Section~\ref{sec:ops}), consistent
    with a defensive promotional response to falling demand.
\end{itemize}

For modelling purposes, we treat the \textbf{post-2018 regime as the relevant
data-generating process} and apply exponentially decaying sample weights to
de-emphasise the pre-break years.

\subsection{Daily Revenue Distribution}

The raw daily revenue distribution is right-skewed (mean $\approx$ 383\,M~VND,
median $\approx$ 286\,M~VND), consistent with the typical log-normal structure of
retail data.
After a $\log$ transformation, the distribution passes a D'Agostino normality
test with $p = 0.31 > 0.05$, confirming approximate log-normality with
$\mu = 19.47$, $\sigma = 0.68$.

This has a direct modelling implication: training on $\log(\text{Revenue})$ and
exponentiating predictions is preferable to modelling raw revenue, as it
naturally enforces non-negativity and makes error distributions more Gaussian.

\subsection{STL Decomposition}

We apply Seasonal-Trend decomposition using LOESS \citep{cleveland1990stl}
with \texttt{period=365} and \texttt{robust=True} to the full daily series.
The variance decomposition is:

\begin{equation}
  \underbrace{y_t}_{\text{observed}} \;=\;
  \underbrace{T_t}_{\text{trend}} \;+\;
  \underbrace{S_t}_{\text{seasonal}} \;+\;
  \underbrace{R_t}_{\text{residual}}
\end{equation}

\begin{table}[H]
\centering
\caption{STL variance decomposition and component strength metrics.}
\label{tab:stl}
\begin{tabular}{lrr}
\toprule
\textbf{Component} & \textbf{\% total variance} & \textbf{Strength (F)}\\
\midrule
Trend ($T_t$)    & 78.3 & $F_T = 1 - \tfrac{\text{Var}(R_t)}{\text{Var}(T_t + R_t)} = 0.941$ \\
Seasonal ($S_t$) & 18.6 & $F_S = 1 - \tfrac{\text{Var}(R_t)}{\text{Var}(S_t + R_t)} = 0.627$ \\
Residual ($R_t$) &  3.1 & — \\
\bottomrule
\end{tabular}
\end{table}

The trend component dominates because of the long upward-then-downward arc.
The seasonal component has a strength of $F_S = 0.627$ — not overwhelming, but
clearly present and forecastable.
Critically, the residual accounts for only 3.1\% of variance, meaning the data
is highly structured and not dominated by noise.
This is good news for forecasting.

\reportfigure{saved_04_stl_decomposition.png}{STL decomposition of daily revenue.}{the separation of long-run trend, annual seasonality, and low residual noise, and supports using decomposition-aware features in the forecasting pipeline.}

\subsection{Autocorrelation Structure}

The autocorrelation function (ACF) of the raw daily series decays slowly,
consistent with a trend-nonstationary process.
Key ACF values:

\begin{table}[H]
\centering
\caption{Selected ACF values — raw $y_t$ and differenced $\Delta\log y_t$.}
\label{tab:acf}
\begin{tabular}{lrr}
\toprule
\textbf{Lag (days)} & \textbf{ACF($y_t$)} & \textbf{ACF($\Delta\log y_t$)}\\
\midrule
1   & 0.874 & 0.031 \\
7   & 0.824 & \textbf{0.142} \\
30  & 0.752 & 0.089 \\
90  & 0.641 & 0.063 \\
180 & 0.553 & 0.047 \\
365 & \textbf{0.737} & \textbf{0.241} \\
548 & 0.662 & 0.198 \\
730 & \textbf{0.648} & \textbf{0.183} \\
\bottomrule
\end{tabular}
\end{table}

Several findings stand out from this table.
First, lag-365 has an \emph{unusually high} ACF of 0.737 in the raw series and
0.241 in the differenced series — far exceeding lag-1 in the differenced case.
This confirms that the dominant signal for forecasting a day 12+ months ahead is
\emph{the same calendar day in the previous year}, not any short-term autocorrelation.
Second, lag-730 (exactly 2 years) also retains strong autocorrelation (0.648/0.183),
suggesting multi-year memory and motivating the use of both 1-year and 2-year lag
features.
Third, the weekly peak at lag-7 in the differenced series confirms a day-of-week
effect, though it is weaker than the annual component.

\reportfigure{saved_05_acf_pacf.png}{Autocorrelation structure of the revenue series.}{strong persistence at annual and biennial horizons in the autocorrelation structure and supports the inclusion of lag-365 and lag-730 features.}

\subsection{Stationarity}

\begin{table}[H]
\centering
\caption{Stationarity tests on $\log(y_t)$ and its first difference.}
\label{tab:stationarity}
\begin{tabular}{llrrrl}
\toprule
\textbf{Series} & \textbf{Test} & \textbf{Statistic} & \textbf{p-value} & \textbf{Critical (5\%)} & \textbf{Conclusion}\\
\midrule
$\log y_t$        & ADF (AIC)  & $-2.41$ & 0.137 & $-2.86$ & Non-stationary \\
$\log y_t$        & KPSS (c+t) & $0.088$ & $>0.10$ & 0.146 & Trend-stationary \\
$\Delta\log y_t$  & ADF (AIC)  & $-38.2$ & $<0.001$ & $-2.86$ & \textcolor{mgreen}{Stationary} \\
\bottomrule
\end{tabular}
\end{table}

The ADF test cannot reject the null of a unit root in $\log y_t$, but KPSS
fails to reject trend-stationarity.
This is the classic ambiguity: the series sits close to the boundary between
integrated-with-drift and trend-stationary.
After one difference, the ADF unambiguously rejects ($p < 0.001$).
In practice, we work with both $\log y_t$ and $\Delta \log y_t$ as targets
depending on the model class.

\subsection{Periodogram — Spectral Dominant Cycles}

The power spectral density of $\log y_t$ shows well-defined peaks at the
following periods:

\begin{table}[H]
\centering
\caption{Top dominant periods from the periodogram of $\log(y_t)$.}
\label{tab:periodogram}
\begin{tabular}{lrl}
\toprule
\textbf{Period (days)} & \textbf{Approximate frequency} & \textbf{Interpretation}\\
\midrule
\textbf{365.0} & Annual & Primary seasonal cycle \\
\textbf{182.5} & Bi-annual & Semi-annual shoulder \\
\textbf{91.3}  & Quarterly & Quarter-end effects \\
\textbf{30.4}  & Monthly & Month-start/end effects \\
\textbf{7.0}   & Weekly & Day-of-week rhythm \\
\textbf{730.0} & Biennial & Even/odd year anomaly \\
\bottomrule
\end{tabular}
\end{table}

The 730-day (biennial) peak is especially interesting — it does not appear in
most retail periodograms and points directly to the August even/odd anomaly
discussed in Section~\ref{subsec:aug}.

\reportfigure{saved_06_periodogram.png}{Dominant spectral cycles in revenue.}{annual, semi-annual, quarterly, weekly, and biennial periodicities in the revenue series and supports frequency-aware feature engineering.}

\subsection{Monthly Seasonal Shape}
\label{subsec:shape}

Normalising each year's monthly revenue by its annual mean reveals that the
\emph{shape} of the seasonal cycle is astonishingly consistent across all 11
years, even through the 2018--2019 structural break.
The profile shows:

\begin{itemize}[leftmargin=1.5em, itemsep=2pt]
  \item \textbf{Peak months}: March--June
    (normalised index $\approx 1.25$--$1.35\times$ the annual mean)
  \item \textbf{Trough months}: November--February
    (normalised index $\approx 0.55$--$0.70\times$ the annual mean)
  \item \textbf{Peak-to-trough ratio}: $\approx 2.7\times$
\end{itemize}

\reportfigure{saved_07_seasonal_shape.png}{Normalised monthly seasonal shape across years.}{that the intra-year seasonal shape remains stable even when the absolute revenue level changes, and supports separating seasonal form from level in the model design.}

\begin{insightbox}[Critical modelling insight — use the shape, not the level]
Because the seasonal \emph{shape} is stable while the seasonal \emph{level} is
not (due to the regime break), the optimal forecasting strategy is to:
(1)~estimate the current level via recent-data lag features,
(2)~apply the stable seasonal shape as a multiplier.
This is precisely what our seasonal lookup tables achieve (Section~\ref{sec:fe}).
\end{insightbox}

\subsection{Day-of-Week Effect}

In the recent regime (2019--2022), day-of-week multipliers relative to the
grand mean are remarkably compressed:

\begin{center}
\small
\begin{tabular}{lccccccc}
\toprule
& Mon & Tue & Wed & Thu & Fri & Sat & Sun \\
\midrule
Multiplier & 0.72 & 0.70 & 0.71 & 0.73 & 0.71 & 0.73 & 0.69 \\
\bottomrule
\end{tabular}
\end{center}

The amplitude of the day-of-week cycle is only $\pm 3$\% around the mean —
far smaller than the $\pm 35$\% annual seasonal amplitude.
While the effect is statistically significant given 3,833 observations, its
practical contribution to forecast error is small.
It is worth including as a feature but should not be over-weighted.

\subsection{The August Biennial Anomaly}
\label{subsec:aug}

The most structurally unusual finding in the dataset is a biennial alternation
in August revenue.
In the recent regime (2019--2022):

\begin{center}
\small
\begin{tabular}{lrr}
\toprule
Year & August Revenue (B VND) & Type \\
\midrule
2019 & 75 & Odd \\
2020 & 117 & Even \\
2021 & 72 & Odd \\
2022 & 120 & Even \\
\bottomrule
\end{tabular}
\end{center}

The Pearson correlation between \texttt{is\_even\_year} (binary) and August
revenue is $r = 0.97$ — essentially deterministic.
Even years (2020, 2022) average $\sim$118~B~VND in August;
odd years (2019, 2021) average $\sim$74~B~VND.
This is a \textbf{59\% difference} that would cause a naive seasonal model to
produce massive errors in August 2023 (odd) and August 2024 (even).

The pattern likely reflects biennial large-scale promotional campaigns
(e.g., end-of-summer clearance events held every two years), corroborated by
the biennial peak visible in the periodogram.
Since 2023 is an odd year, the model should predict a \emph{below-average}
August 2023 and an \emph{above-average} August 2024.

\reportfigure{saved_08_dow_august_biennial.png}{The August biennial anomaly.}{the pronounced even-year versus odd-year August divergence and supports explicit biennial anomaly flags in the forecasting feature set.}

% ─────────────────────────────────────────────────────────────────────────────
\section{Customer Analysis}
\label{sec:customers}
% ─────────────────────────────────────────────────────────────────────────────

\subsection{Acquisition and Cumulative Growth}

New customer acquisition grew steadily from 2012 to 2017, peaking at
$\approx$700--800 new signups per month.
Acquisition then declined sharply in 2018--2019 (consistent with the revenue
break) before stabilising at $\approx$300--400/month in the recent regime.
The cumulative customer base still shows upward momentum, reaching
$\sim$50,000 registered customers by end of 2022.

\subsection{Purchase Frequency Distribution}

The distribution of orders per customer is heavily right-skewed, consistent
with the empirical power-law typical of e-commerce:

\begin{itemize}[leftmargin=1.5em, itemsep=2pt]
  \item \textbf{One-time buyers}: $\approx 45\%$ of all registered customers
    placed exactly one order. This is high by industry standards and suggests
    significant friction in the second-purchase journey.
  \item \textbf{Median}: 2 orders per customer.
  \item \textbf{Mean}: $\approx$3.8 orders per customer (pulled right by power users).
  \item \textbf{Top 20\% of customers}: account for $\approx$65\% of all orders
    (Lorenz curve / Gini $\approx 0.52$).
\end{itemize}

\subsection{Customer Demographics}

Customer demographics are broadly split:
\textit{acquisition channel} shows a roughly even split between organic search,
social media referral, and direct/email;
\textit{age group} skews Gen-Z (18--24) and Millennial (25--34);
\textit{gender} is approximately 60/40 female/male;
\textit{loyalty tier} shows that approximately 30\% of customers have reached
Silver or above — a relatively high tier penetration indicating an engaged core.

\subsection{RFM Segmentation}

Applying standard RFM (Recency-Frequency-Monetary) quintile scoring yields five
segments (Table~\ref{tab:rfm}).

\begin{table}[H]
\centering
\caption{RFM segment summary. Champions and Loyal customers together represent
38\% of the base but account for over 70\% of revenue.}
\label{tab:rfm}
\small
\begin{tabular}{lrrrr}
\toprule
\textbf{Segment} & \textbf{N customers} & \textbf{Avg Recency (d)}
  & \textbf{Avg Orders} & \textbf{Total Rev (\%)}\\
\midrule
Champions & $\sim$4,200  & 38  & 12.4 & 42.1 \\
Loyal     & $\sim$14,800 & 112 & 5.8  & 29.4 \\
Recent    & $\sim$8,600  & 21  & 1.4  & 9.3  \\
At Risk   & $\sim$11,400 & 345 & 4.1  & 14.7 \\
Lost      & $\sim$11,000 & 641 & 1.2  & 4.5  \\
\bottomrule
\end{tabular}
\end{table}

The large ``Lost'' and ``At Risk'' cohorts (together $\sim$45\% of customers)
are a latent revenue opportunity.
Win-back campaigns targeting these segments may partially explain the revenue
recovery observed in 2022.

\reportfigure{saved_13_rfm_segments.png}{RFM segment concentration.}{that Champions and Loyal customers account for a disproportionate share of revenue and supports segment-aware retention strategy and feature design.}

% ─────────────────────────────────────────────────────────────────────────────
\section{Product and Category Analysis}
\label{sec:products}
% ─────────────────────────────────────────────────────────────────────────────

\subsection{Category Revenue and Margin}

The product catalogue is organised into four major categories.
Revenue concentration is extreme:

\begin{table}[H]
\centering
\caption{Revenue, order count, and estimated gross margin by product category.}
\label{tab:categories}
\begin{tabular}{lrrrr}
\toprule
\textbf{Category} & \textbf{Revenue (\%)} & \textbf{Orders (\%)}
  & \textbf{Unique SKUs} & \textbf{Est. Margin \%}\\
\midrule
Streetwear   & 78.1 & 72.4 & $\sim$620 & 9.3 \\
Outdoor      & 10.8 & 13.7 & $\sim$180 & 18.4 \\
GenZ Fashion & 7.6  & 9.2  & $\sim$140 & 22.1 \\
Kids         & 3.5  & 4.7  & $\sim$80  & 15.7 \\
\bottomrule
\end{tabular}
\end{table}

The business is overwhelmingly Streetwear-dependent (78\% of revenue) yet this
category carries the \emph{lowest} gross margin at $\approx$9.3\%.
Outdoor and GenZ Fashion, by contrast, generate more than double the margin.
A product mix shift toward higher-margin categories would significantly improve
profitability without necessarily growing top-line revenue.

\reportfigure{saved_14_category_revenue.png}{Category revenue and margin mix.}{category-level revenue concentration and margin heterogeneity, and supports the conclusion that Streetwear dominates the business mix.}

\subsection{Price and COGS Distribution}

Product prices follow an approximately log-normal distribution with a median of
$\approx$320,000~VND and a heavy right tail extending to luxury items above
2,000,000~VND.
COGS tracks price linearly with a mean product-level margin of $\approx$14.8\%,
consistent with the aggregate margin observed in the sales table — a useful
cross-validation of data consistency.

% ─────────────────────────────────────────────────────────────────────────────
\section{Operational Signals}
\label{sec:ops}
% ─────────────────────────────────────────────────────────────────────────────

\subsection{Promotions}
\label{subsec:promos}

Approximately 80 distinct promotional campaigns were run over the decade.
Three dominant \texttt{promo\_type} categories account for $>$90\% of
campaigns: percentage discount, fixed-amount coupon, and free shipping.
Campaign durations range from 3 to 45 days, with a median of 14 days.

Promo penetration (fraction of order-item lines carrying a promo code)
alternates biannually:

\begin{center}
\small
\begin{tabular}{lccccccccccc}
\toprule
Year & 2012 & 2013 & 2014 & 2015 & 2016 & 2017 & 2018 & 2019 & 2020 & 2021 & 2022 \\
\midrule
Promo \% & 28.3 & 31.2 & 33.1 & 29.7 & 25.4 & 27.8 & 33.6 & 46.1 & 38.2 & 48.7 & 39.4 \\
\bottomrule
\end{tabular}
\end{center}

The sharp increase in promo penetration starting 2019 (from $\sim$33\% to 46--49\%)
is directly correlated with the revenue regime break — suggesting that promotions
were deployed defensively to cushion the demand decline rather than proactively
to grow revenue.
Importantly, \emph{promo penetration is negatively correlated with revenue}
(Pearson $r \approx -0.61$) in the full series, cautioning against naive
conclusions that more promotions drive more sales.

\begin{insightbox}[Promotions as a feature — not a cause]
The data does not support the view that promos \emph{cause} revenue growth.
The correlation is confounded by the regime break.
In the post-break regime, promotions appear to provide a floor rather than
a ceiling for revenue.
For forecasting, \texttt{is\_even\_year} (which captures the biennial promo
intensity pattern) is more parsimonious and equally predictive.
\end{insightbox}

\reportfigure{saved_18_promo_penetration.png}{Promotion penetration over time.}{the rise in promotion penetration after 2019 alongside the lower-revenue regime, and supports treating promotions as regime-correlated rather than purely exogenous signals.}

\subsection{Shipping Performance}

Shipping days (ship date $\to$ delivery date) follow a roughly discrete
distribution:

\begin{itemize}[leftmargin=1.5em, itemsep=2pt]
  \item Median shipping time: \textbf{3 days}
  \item 75th percentile: 5 days
  \item 95th percentile: 8 days
  \item Out-of-range ($>$30 days): $<$0.5\% of shipments (likely data errors,
    filtered)
\end{itemize}

Mean shipping days have improved slightly over the decade (from $\sim$4.2~days
in 2012 to $\sim$3.1~days in 2022), consistent with investment in logistics
infrastructure.

\subsection{Return Rate and Reasons}

The aggregate monthly return rate averages \textbf{5.8\%} (returns / orders)
with relatively low variance over time.
A modest increase in return rates is visible from 2020 onward, possibly related
to higher discount-driven purchases of less-deliberate items.

The top five return reasons are:

\begin{enumerate}[leftmargin=1.5em, itemsep=2pt]
  \item Wrong size / poor fit (35\%)
  \item Product quality below expectation (22\%)
  \item Wrong item received (14\%)
  \item Changed mind / no longer needed (18\%)
  \item Delivery damage (11\%)
\end{enumerate}

\begin{warningbox}[Operational concern — size returns]
Wrong size is the \#1 return reason at 35\%.
Yet return rates are uniform across all size categories, suggesting that the
issue is \emph{systemic catalogue inaccuracy} (misleading size charts, model
photos not representative) rather than manufacturing defects.
This is addressable via UX improvements and would directly reduce operational
costs.
\end{warningbox}

\subsection{Customer Reviews and Ratings}

The rating distribution is left-skewed (positivity bias, common in fashion
retail):

\begin{center}
\small
\begin{tabular}{lrrrrr}
\toprule
Rating & 1★ & 2★ & 3★ & 4★ & 5★ \\
\midrule
Fraction & 3.2\% & 5.1\% & 12.8\% & 38.6\% & 40.3\% \\
\bottomrule
\end{tabular}
\end{center}

Overall mean rating is $4.11$ out of 5.
Monthly average ratings are stable over time with no secular trend, indicating
consistent product quality.

\subsection{Web Traffic and Conversion}

Web traffic (sessions) grew roughly in parallel with revenue through 2016,
then declined.
The monthly order conversion rate (orders / sessions) has been relatively stable
at $\approx$\textbf{3.2--4.8\%}, declining slightly in the post-break regime.
This suggests that the revenue decline in 2019 was driven by reduced traffic
rather than worsening conversion — a demand acquisition problem, not a
funnel efficiency problem.

\subsection{COD Cancellation Rate}
\label{subsec:cod}

Cash-on-Delivery (COD) exhibits a cancellation rate of \textbf{$\approx$16--18\%},
roughly twice that of card or e-wallet payments ($\approx$7--9\%).
This is a known characteristic of Vietnamese and Southeast Asian markets where
COD remains the dominant payment method for lower-income segments.

\begin{table}[H]
\centering
\caption{Cancellation rate by payment method (recent regime, 2019--2022).}
\label{tab:cancel}
\begin{tabular}{lrr}
\toprule
\textbf{Payment method} & \textbf{Cancellation rate \%} & \textbf{Relative to mean}\\
\midrule
Cash-on-Delivery (COD) & 17.3 & \textcolor{mred}{$+2.1\times$} \\
Bank transfer          & 8.4  & 1.0$\times$ \\
Credit / debit card    & 7.6  & $0.92\times$ \\
E-wallet (MoMo, etc.)  & 6.8  & $0.82\times$ \\
\bottomrule
\end{tabular}
\end{table}

The revenue impact of COD cancellations is material: assuming COD represents
$\sim$55\% of orders and 17\% of those cancel, approximately 9\% of
\emph{gross} orders never generate revenue.
This drag is already embedded in the \texttt{sales.Revenue} figures (cancelled
orders are not counted), but it creates a noisy relationship between order
\emph{placement} signals and \emph{realised} revenue.

\reportfigure{saved_21_cod_cancellation.png}{Cancellation rates by payment method.}{materially higher cancellation rates for COD orders than for card or e-wallet payments and supports payment-method-aware operational risk controls.}

% ─────────────────────────────────────────────────────────────────────────────
\section{Preprocessing and Feature Engineering}
\label{sec:fe}
% ─────────────────────────────────────────────────────────────────────────────

This section documents every feature engineering decision, the rationale behind
it, and the leakage implications for the 548-day forecast horizon.
The key constraint is that \emph{no information from after a given prediction
date may enter that day's feature vector}.
Since the forecast starts 2023-01-01, any feature that references data within
12 months of that date must use lag-365 or greater.

\subsection{Calendar Features}

\paragraph{Raw calendar.}
Year, month (1--12), day-of-month (1--31), day-of-week (0=Monday, 6=Sunday),
day-of-year (1--365), week-of-year (1--52), quarter (1--4).

\paragraph{Binary boundary flags.}
The binary calendar indicators are \texttt{is\_month\_end},
\texttt{is\_month\_start}, \texttt{is\_quarter\_end},
\texttt{is\_quarter\_start}, \texttt{is\_year\_end},
\texttt{is\_year\_start}, and \texttt{is\_weekend}.
These capture the well-documented month-end and quarter-end spikes common
in fashion retail (payday effects, reporting period effects).

\paragraph{Cyclical encodings.}
To avoid treating month 12 as numerically ``far from'' month 1, all periodic
calendar variables are encoded as:
\begin{equation}
  \text{col\_sin} = \sin\!\left(\frac{2\pi \cdot \text{col}}{\text{period}}\right),
  \qquad
  \text{col\_cos} = \cos\!\left(\frac{2\pi \cdot \text{col}}{\text{period}}\right)
\end{equation}
Applied to month ($T=12$), day-of-week ($T=7$), week-of-year ($T=52$), and
day-of-year ($T=365$), yielding 8 cyclical features.

\subsection{Vietnamese Holiday and Tet Features}

Vietnamese public holidays (New Year Jan~1, Reunification Apr~30, Labour Day
May~1, National Day Sep~2) are encoded as a binary flag
\texttt{is\_vn\_holiday}.
The Lunar New Year (\textit{Tet Nguyen Dan}) is the single most important
demand driver in Vietnamese retail.
We encode proximity to Tet via:

\begin{equation}
  \texttt{days\_to\_tet}_t = t - \hat{\tau}_{\text{Tet}(year(t))}
\end{equation}

where $\hat{\tau}_{\text{Tet}(y)}$ is the Gregorian date of the first day of
Tet in year $y$ (approximated from astronomical tables through 2024).
Negative values indicate pre-Tet buildup; positive values indicate post-Tet
recovery.
Binary thresholds \texttt{is\_tet\_week} ($|\texttt{days\_to\_tet}| \le 7$)
and \texttt{is\_tet\_month} ($-30 \le \texttt{days\_to\_tet} \le 30$) are
derived features.

\subsection{Regime and Parity Flags}

\begin{itemize}[leftmargin=1.5em, itemsep=2pt]
  \item \texttt{is\_recent\_regime}: 1 if year $\ge$ 2019, else 0
  \item \texttt{is\_even\_year}: 1 if year is even (2020, 2022, 2024 \ldots), else 0
  \item \texttt{is\_august}: 1 if month == 8
  \item \texttt{is\_aug\_even}: interaction flag (August $\times$ even year)
  \item \texttt{is\_aug\_odd}: interaction flag (August $\times$ odd year)
\end{itemize}

These four flags collectively encode the biennial August anomaly
(Section~\ref{subsec:aug}) in a form any tree or linear model can exploit.

\subsection{Long-Horizon Lag Features}
\label{subsec:lags}

Since the forecast horizon begins more than 365 days after the last training
observation, \emph{only lags $\ge$ 365 are leakage-safe for all forecast dates}.
We include:

\begin{table}[H]
\centering
\caption{Long-horizon lag features. All are computed on both raw Revenue and
$\log(\text{Revenue})$.}
\label{tab:lags}
\begin{tabular}{lll}
\toprule
\textbf{Feature} & \textbf{Lag (days)} & \textbf{Rationale}\\
\midrule
\texttt{rev\_lag\_365}     & 365 & Same calendar date, 1 year prior \\
\texttt{rev\_lag\_366}     & 366 & Leap-year correction \\
\texttt{rev\_lag\_364}     & 364 & Same weekday, 52 weeks prior \\
\texttt{rev\_lag\_548}     & 548 & 18 months prior \\
\texttt{rev\_lag\_730}     & 730 & Same date, 2 years prior \\
\texttt{rev\_roll7\_lag365} & 365+7d window & Smoothed 1-week centred on lag-365 \\
\texttt{rev\_roll30\_lag365}& 365+30d window & Smoothed 1-month centred on lag-365 \\
\texttt{rev\_roll90\_lag365}& 365+90d window & Smoothed 1-quarter centred on lag-365 \\
\texttt{rev\_lag\_365\_730\_mean} & avg(365, 730) & Blended 1+2 year signal \\
\bottomrule
\end{tabular}
\end{table}

The blended lag (\texttt{rev\_lag\_365\_730\_mean}) is particularly powerful
because it partially averages out the biennial August anomaly — if 2022 was an
even-year August and 2021 was an odd-year August, blending them provides a
``neutral'' baseline from which the \texttt{is\_aug\_even} flag can then deviate.

\subsection{Seasonal Lookup Tables}

For each of four grouping keys, we compute mean revenue from the recent regime
(2019--2022) and apply a level-correction factor $\lambda$ to account for the
upward trend within that regime:

\begin{equation}
  \lambda = \frac{\bar{y}_{[-365,0]}}{\bar{y}_{\text{2019-2022}}}
\end{equation}

The four lookup tables are:

\begin{enumerate}[leftmargin=1.5em, itemsep=2pt]
  \item \textbf{Day-of-year mean} (\texttt{seasonal\_doy\_mean}): 365 values,
    captures the full annual shape
  \item \textbf{Month $\times$ Day-of-week mean} (\texttt{seasonal\_month\_dow\_mean}):
    $12 \times 7 = 84$ cells, captures month-within-year and DoW jointly
  \item \textbf{Month $\times$ Year-parity mean} (\texttt{seasonal\_month\_parity\_mean}):
    $12 \times 2 = 24$ cells, directly encodes the biennial pattern by month
  \item \textbf{Week-of-year mean} (\texttt{seasonal\_week\_mean}): 52 values,
    coarser alternative to DoY
\end{enumerate}

\begin{insightbox}[Level correction is critical]
Without level correction, the lookup tables estimate what revenue
\emph{was} in 2019--2022, not what it \emph{will be} in 2023--2024.
By scaling by $\lambda$, the lookups inherit the most recent 12-month level,
making them forward-looking baseline estimates rather than historical averages.
\end{insightbox}

\subsection{Sample Weights}

To emphasise the recent regime without completely discarding the historical
data (which contains valuable seasonal shape information), we apply
exponentially decaying sample weights with a two-year half-life:

\begin{equation}
  w_t = \exp\!\left(\frac{-0.5 \cdot (t_{\max} - t)}{730}\right)
  \times
  \begin{cases} 0.3 & \text{if } year(t) < 2019 \\ 1.0 & \text{otherwise} \end{cases}
\end{equation}

This ensures that December 2022 receives weight $w = 1.0$, December 2020
receives $w \approx 0.60$, and December 2018 receives $w \approx 0.30 \times
0.58 \approx 0.17$, making pre-break data nearly inconsequential without
deleting it entirely.

\reportfigure{saved_22_sample_weights.png}{Regime-aware sample weighting schedule.}{the weighting schedule that downweights older observations, especially the pre-2019 regime, and supports more relevant model fitting for the forecast horizon.}

\subsection{Walk-Forward Cross-Validation}

Because the test period is both long (548 days) and entirely beyond the training
data, standard $k$-fold cross-validation is inappropriate — it would permit
future information to leak into the past.
We instead use the following walk-forward scheme:

\begin{table}[H]
\centering
\caption{Walk-forward CV splits. The validation window in each fold approximates
the length and statistical properties of the actual test period.}
\label{tab:cv}
\begin{tabular}{lllll}
\toprule
\textbf{Fold} & \textbf{Train start} & \textbf{Train end}
  & \textbf{Val start} & \textbf{Val end}\\
\midrule
Fold 1 & 2012-07-04 & 2020-06-30 & 2020-07-01 & 2021-12-31 \\
Fold 2 & 2012-07-04 & 2021-06-30 & 2021-07-01 & 2022-12-31 \\
Fold 3 & 2012-07-04 & 2021-12-31 & 2022-01-01 & 2022-12-31 \\
\midrule
\textbf{Test} & 2012-07-04 & 2022-12-31 & 2023-01-01 & 2024-07-01 \\
\bottomrule
\end{tabular}
\end{table}

Fold~3 is the most important proxy for the actual test performance: both the
training window and the validation window closely match the actual
train/test split, and the validation period (Jan--Dec 2022) includes the same
seasonal challenges as the first year of the forecast (Jan--Dec 2023).

\reportfigure{saved_23_cv_splits.png}{Walk-forward cross-validation design.}{the leakage-safe walk-forward validation design and supports honest estimation of out-of-sample forecasting performance.}

% ─────────────────────────────────────────────────────────────────────────────
\section{Consolidated Key Findings}
\label{sec:findings}
% ─────────────────────────────────────────────────────────────────────────────

We distil the EDA into eight falsifiable hypotheses that directly guide model
design:

\begin{enumerate}[leftmargin=1.5em, itemsep=6pt]

\item[\textbf{H1}] \textbf{Regime-aware weighting will outperform
  uniform weighting.}
  Pre-2019 data is drawn from a different level regime.
  Any model trained without downweighting these years will underpredict 2023--2024
  revenue.

\item[\textbf{H2}] \textbf{Lag-365 is the single strongest predictive feature.}
  The annual ACF of 0.737 (raw) and the stable seasonal shape both confirm that
  ``what happened on this calendar day last year'' is the most informative
  single signal.

\item[\textbf{H3}] \textbf{Lag-730 adds independent information over Lag-365.}
  The two-year ACF of 0.648 is not simply a consequence of chaining two
  one-year autocorrelations.
  Including both lags and their average will improve point estimates.

\item[\textbf{H4}] \textbf{The August biennial flag must be explicitly encoded.}
  Failure to include \texttt{is\_aug\_even}/\texttt{is\_aug\_odd} will result in
  a $\sim$59\% forecasting error in August months, which will dominate overall
  MAPE.

\item[\textbf{H5}] \textbf{Tet proximity is the highest-impact holiday feature.}
  Vietnamese Lunar New Year is the dominant cultural shopping event.
  Generic ``is-public-holiday'' encoding will miss the 2--3 week buildup and the
  post-Tet trough.

\item[\textbf{H6}] \textbf{The seasonal shape is stable; the level is not.}
  Models that learn the shape (e.g., via normalised seasonal lookups) and
  separately learn the level (e.g., via recent lags) will generalise better than
  models that try to directly predict absolute revenue.

\item[\textbf{H7}] \textbf{Promo flags are weak exogenous predictors.}
  Promo penetration is endogenous to the demand regime and correlated with, not
  causal of, revenue.
  Binary promo flags may add marginal value but \texttt{is\_even\_year}
  (capturing the biennial promo cycle) is a more parsimonious predictor.

\item[\textbf{H8}] \textbf{Log-target training is preferable to raw-target.}
  The log-normal distribution of daily revenue, combined with the
  heteroscedasticity visible in residuals (variance scales with level), means
  that training on $\log(\text{Revenue})$ and exponentiating predictions will
  produce more calibrated forecasts with lower MAPE.

\end{enumerate}

% ─────────────────────────────────────────────────────────────────────────────
\section{Feature Engineering Summary}
\label{sec:summary}
% ─────────────────────────────────────────────────────────────────────────────

Table~\ref{tab:features} provides a complete inventory of the final feature
matrix.

\reportfigure{saved_27_feature_correlations.png}{Feature-target correlation summary.}{which engineered features are most strongly associated with the target and supports prioritising seasonal, lag-based, and regime-aware predictors.}

\begin{table}[H]
\centering
\caption{Final feature matrix — complete inventory.
All features are leakage-safe for the 2023-01-01 → 2024-07-01 forecast horizon.}
\label{tab:features}
\small
\begin{tabular}{>{\raggedright\arraybackslash}p{4.4cm}>{\raggedright\arraybackslash}p{1.7cm}>{\raggedright\arraybackslash}p{6.9cm}}
\toprule
\textbf{Feature group} & \textbf{Count} & \textbf{Description}\\
\midrule
Calendar (raw)          & 9  & year, month, dom, dow, doy, woy, quarter, etc. \\
Calendar (binary)       & 8  & is\_weekend, is\_month\_end, is\_month\_start, etc. \\
Cyclical encodings      & 8  & sin/cos for month, dow, woy, doy \\
Regime flags            & 5  & is\_recent\_regime, is\_even\_year, is\_august, etc. \\
Tet proximity           & 4  & days\_to\_tet, is\_tet\_week, is\_tet\_month, countdown \\
VN public holidays      & 1  & is\_vn\_holiday \\
Long-horizon lags (Rev) & 5  & lag 364/365/366/548/730 \\
Long-horizon lags (log) & 5  & log-transformed versions of above \\
Rolling means (Rev)     & 12 & roll7/14/30/90 $\times$ lag365/548/730 \\
Rolling means (log)     & 12 & log-transformed versions of above \\
Blended lags            & 2  & avg(lag365, lag730) raw + log \\
Seasonal lookup (Rev)   & 4  & doy, month×dow, month×parity, week \\
Seasonal lookup (log)   & 4  & log versions of above \\
Sample weight           & 1  & exp decay, old regime ×0.3 \\
\midrule
\textbf{Total}          & \textbf{80} & \\
\bottomrule
\end{tabular}
\end{table}

\paragraph{Top features by correlation with $\log(\text{Revenue})$}
(computed on the full training set):

\begin{enumerate}[leftmargin=1.5em, itemsep=2pt]
  \item \texttt{seasonal\_month\_parity\_mean} ($r = 0.74$)
  \item \texttt{log\_rev\_lag\_365} ($r = 0.71$)
  \item \texttt{seasonal\_doy\_mean} ($r = 0.68$)
  \item \texttt{log\_rev\_roll30\_lag365} ($r = 0.67$)
  \item \texttt{seasonal\_month\_dow\_mean} ($r = 0.64$)
  \item \texttt{is\_aug\_even} ($r = 0.41$)
  \item \texttt{is\_aug\_odd} ($r = -0.38$)
  \item \texttt{is\_tet\_month} ($r = -0.31$)
  \item \texttt{month\_sin} ($r = 0.29$)
  \item \texttt{is\_even\_year} ($r = 0.27$)
\end{enumerate}

The top three features all capture the annual seasonal pattern through different
lenses (monthly parity lookup, direct lag, rolling smoothed lag), confirming
that the seasonal signal is the dominant forecastable component.

% ─────────────────────────────────────────────────────────────────────────────
\section{Conclusion}
\label{sec:conclusion}
% ─────────────────────────────────────────────────────────────────────────────

This EDA has transformed a decade of raw transactional records into a
well-understood dataset with clear modelling implications.
The core insight is that Vietnamese fashion e-commerce revenue is \emph{not}
a noisy random walk — it contains rich, stable, multi-scale periodic structure
embedded within a single regime-breaking trend.
A forecasting model that honours this structure (regime-aware weighting, annual
lags, biennial parity flags, Tet proximity, and level-corrected seasonal
lookups) should achieve substantially lower error than off-the-shelf time-series
methods that ignore the domain context.

The feature matrix produced here — 80 leakage-safe features across 3,833
training observations — is ready for downstream gradient boosting, ARIMA-X,
neural forecasting, or ensemble methods.
The walk-forward cross-validation scheme provides an honest performance estimate
and guards against both overfitting and optimism bias.

\paragraph{Next steps.}
\begin{itemize}[leftmargin=1.5em, itemsep=2pt]
  \item Train and tune a LightGBM baseline on the 80-feature matrix with
    sample weights and log-target.
  \item Blend with a classical Fourier-extended regression for interpretability.
  \item Evaluate post-processing: does applying the seasonal shape to the model's
    trend estimate outperform the model's native seasonal prediction?
  \item Explicitly model the August biennial anomaly as a regime-switching
    component in a hybrid model.
\end{itemize}

% ─────────────────────────────────────────────────────────────────────────────
% Bibliography
% ─────────────────────────────────────────────────────────────────────────────
\bibliographystyle{abbrvnat}
\begin{thebibliography}{9}

\bibitem{cleveland1990stl}
R.~B. Cleveland, W.~S. Cleveland, J.~E. McRae, and I.~Terpenning.
\newblock {STL}: A seasonal-trend decomposition procedure based on {LOESS}.
\newblock \emph{Journal of Official Statistics}, 6(1):3--73, 1990.

\bibitem{hyndman2018forecasting}
R.~J. Hyndman and G.~Athanasopoulos.
\newblock \emph{Forecasting: Principles and Practice}.
\newblock OTexts, 2nd edition, 2018.
\newblock Available at \url{https://otexts.com/fpp2/}.

\bibitem{ke2017lightgbm}
G.~Ke, Q.~Meng, T.~Finley, T.~Wang, W.~Chen, W.~Ma, Q.~Ye, and T.-Y. Liu.
\newblock {LightGBM}: A highly efficient gradient boosting decision tree.
\newblock In \emph{Advances in Neural Information Processing Systems (NeurIPS)},
  pages 3146--3154, 2017.

\bibitem{rfm2000}
P.~D. Berger and N.~I. Nasr.
\newblock Customer lifetime value: Marketing models and applications.
\newblock \emph{Journal of Interactive Marketing}, 12(1):17--30, 1998.

\bibitem{makridakis2018m4}
S.~Makridakis, E.~Spiliotis, and V.~Assimakopoulos.
\newblock The {M4} competition: Results, findings, conclusion and way forward.
\newblock \emph{International Journal of Forecasting}, 34(4):802--808, 2018.

\end{thebibliography}

% ─────────────────────────────────────────────────────────────────────────────
\appendix
\section{Stationarity Test Details}
\label{app:stationarity}

\paragraph{ADF test.}
The Augmented Dickey-Fuller test \citep{hyndman2018forecasting} evaluates $H_0$:
unit root present against $H_1$: stationary.
Lag selection by AIC.
On $\log y_t$: statistic $= -2.41$, $p = 0.137$ (cannot reject $H_0$).
On $\Delta \log y_t$: statistic $= -38.2$, $p < 0.001$ (reject $H_0$,
series is integrated of order 1).

\paragraph{KPSS test.}
The Kwiatkowski-Phillips-Schmidt-Shin test evaluates $H_0$: trend-stationary
against $H_1$: unit root.
With \texttt{regression='ct'} (constant + trend): statistic $= 0.088$,
$p > 0.10$ (fail to reject $H_0$, trend-stationary).
The apparent contradiction between ADF and KPSS reflects a series that is
\emph{near} the boundary: long-memory near-integrated, or I(1) with large
MA root.
Both characterisations motivate including lag-365 and lag-730 features that
are insensitive to this distinction.

\section{Tet Date Table}
\label{app:tet}

\begin{center}
\small
\begin{tabular}{ll}
\toprule
\textbf{Year} & \textbf{Tet (first day, Gregorian)}\\
\midrule
2012 & January 23 \\
2013 & February 10 \\
2014 & January 31 \\
2015 & February 19 \\
2016 & February 8 \\
2017 & January 28 \\
2018 & February 16 \\
2019 & February 5 \\
2020 & January 25 \\
2021 & February 12 \\
2022 & February 1 \\
2023 & January 22 \\
2024 & February 10 \\
\bottomrule
\end{tabular}
\end{center}

\end{document}